% Intro + related work, revised 2026-09-11 after reviewer comments (C013-C032), with the ~500-char cut applied (v4: Geoffrey's tracked-change text restored verbatim).
% Abstract unchanged from the current Overleaf version.

Autonomous cyber-physical systems such as cars, drones and industrial robots run a processing pipeline that typically consists of sensing, perception, planning and control. The processing chain from a sensor reading to the resulting actuation must complete within an end-to-end deadline for the system to be functional and safe, and some of its steps offload work to GPUs. Such systems are built on a middleware, a layer between application and OS that provides communication and composition and raises the level of abstraction at which components are written, so that the hundreds of components of the pipeline can be built with reasonable effort.

ROS~2~\cite{macenski-2022-ros2} is one such middleware. It abstracts an application into three levels: the callback is the unit of computation; a node groups the timers, subscriptions, services, clients and other waitables that trigger a component's callbacks, and the publishers they send through; and an executor dispatches the ready callbacks of the nodes assigned to it on one or more threads of an OS process. Three parties decide when each callback runs and on which thread, core or GPU engine, and none of them sees the other two: the OS schedules the executor threads, the executor picks the next ready callback, and the GPU run-time orders the work that callback submits.

Timing analysis of ROS~2 executors is a recent but fast-evolving field, e.g.,~\cite{casini-19-reservationbased,arafat-2026-mtexecutor-cbgroups,teper-2024-end2endmtexecutoranalysisoptimization}. These analyses share one prerequisite: a model of the application. The chains, the callback-to-executor mapping and the callback-group memberships are inputs, so each analysis describes the model it is given, not the application that actually runs.

% NOTE: the clause ", and the ROS~2 layer only in part: never waitables, ... publish/subscribe." near the end of this paragraph is optional: drop it for −156 chars (contribution II already states the coverage).
A complementary line of work focuses on building the application model required by those analyses from a running system. For instance, ROS-Llama~\cite{blass-2021-rosllama-paper} and Abaza et al.~\cite{abaza-2024-ros2modelsynthesis} instrument the ROS~2 libraries and recover the callback graph with measured execution times; LiME\textsuperscript{IT}~\cite{brandenburg-2026-limeit} infers a periodic/sporadic model of ROS~2 callbacks from OS events and probes, without using ROS~2-specific tools. All three consider only a single-threaded executor, which leaves out the thread count and callback-group memberships that a multi-threaded executor adds---and Autoware relies on---to decide which callbacks may run in parallel. A further family of tools builds on the instrumentation provided by \texttt{ros2\_tracing}~\cite{bedard-2022-ros2tracing}: CARET~\cite{kuboichi-2022-caret} records nodes, callbacks and topics and follows individual messages through them, and an extension adds executors and callback groups~\cite{peng-2022-caretcomp}; both leave the chain to the user and report latencies rather than a model. Each of these tools therefore covers at most two of the three layers---OS, ROS~2 and GPU---that a deployed application spans, and the ROS~2 layer only in part: never waitables, which carry intra-process subscriptions and action servers, and no interaction beyond publish/subscribe. None records the accelerator a callback uses, which leaves the time a callback spends waiting for its GPU work indistinguishable from computation; making accelerators first-class in the analysis of ROS~2 systems is an open problem~\cite{yano-2026-rightanalyticalmodelsoftware}.

In this work, we fill that gap by presenting a solution to automatically and systematically extract an expressive model of a ROS~2 application from runtime observations. The model is based on actual execution traces. Therefore, it does not require expert knowledge about the application that needs to be modeled, nor does it rely on potentially incomplete or outdated documentation.

\textbf{Our contributions are:} \textbf{(I)}~a hierarchical model of a deployed ROS~2 application spanning the OS, ROS~2 and GPU layers; \textbf{(II)}~instrumentation that extends \texttt{ros2\_tracing} with tracepoints for object lifetimes, all five kinds of callback an executor dispatches and intra-process delivery, for rclcpp and rclpy alike; \textbf{(III)}~a trace-acquisition and model-generation method that turns the trace into the model, verified against the objects the traced binaries create.
